\documentclass[11pt,a4paper]{article}

\usepackage[T1]{fontenc}
\usepackage{amsmath,amssymb,mathtools}
\usepackage{array}
\usepackage{booktabs}
\usepackage{enumitem}
\usepackage{geometry}
\usepackage{microtype}
\usepackage{xcolor}
\usepackage{hyperref}

\geometry{margin=27mm}
\hypersetup{
  colorlinks=true,
  linkcolor=blue!55!black,
  urlcolor=blue!55!black,
  citecolor=blue!55!black,
  pdftitle={How a Permutation FFT Tames the Gluon Colour Sum},
  pdfauthor={MultipletRecursion tutorial note}
}

\setlist[itemize]{leftmargin=2em}
\setlist[enumerate]{leftmargin=2em}
\newcommand{\ii}{\mathrm{i}}
\newcommand{\Tr}{\operatorname{Tr}}
\newcommand{\cA}{\mathcal{A}}
\newcommand{\cC}{\mathcal{C}}
\newcommand{\cM}{\mathcal{M}}
\newcommand{\code}[1]{\texttt{#1}}
\newenvironment{takeaway}
  {\par\medskip\noindent\begin{minipage}{\linewidth}
   \color{blue!50!black}\textbf{Takeaway.}\ \color{black}}
  {\end{minipage}\par\medskip}

\title{How a Permutation FFT Tames the Gluon Colour Sum\\
\large A beginner's guide to the trace and adjoint implementations}
\author{A tutorial note for the \textsc{AmpliGluon} codes}
\date{\today}

\begin{document}

\maketitle

\begin{abstract}
Full-colour gluon calculations written in a trace or Del Duca--Dixon--Maltoni
(DDM) basis contain a dense colour Gram matrix.  If there are $P$ colour
orders, storing this matrix and applying it directly both appear to require
order $P^2$ resources.  The implementations in \textsc{AmpliGluonTrace} and
\textsc{AmpliGluonAdjoint} avoid that dense calculation by noticing that a
Gram entry depends only on the \emph{relative permutation} between two
orders.  The matrix is therefore a convolution on a symmetric group.  A
non-abelian fast Fourier transform reorganizes that convolution into
independent blocks labelled by Young diagrams.

This note develops the idea from scratch.  It explains what is transformed,
why the Fourier coefficients are matrices, how the colour contraction is
recovered exactly, how the recursive transform is made fast, and how the
trace and adjoint codes differ.  Two examples can be checked by hand.  We
also separate the precise result---linear colour \emph{storage} and
subquadratic structured contraction---from the useful but imprecise slogan
that the cost changes from ``factorial squared'' to ``about factorial.''
\end{abstract}

\tableofcontents

\section{Executive answer}

The short answer is yes: the symmetric-group FFT is the important device
that prevents the trace and adjoint codes from constructing and multiplying
the full dense matrix in
\begin{equation}
  Q = \cA^\dagger G\cA.
  \label{eq:headline}
\end{equation}
However, there are two logically separate ingredients:

\begin{enumerate}
  \item \textbf{Relative-permutation symmetry} shows that all $P^2$
        entries of $G$ are determined by only $P$ numbers.
  \item \textbf{The symmetric-group FFT} reorganizes the multiplication by
        $G$ into independent, much smaller representation blocks.
\end{enumerate}

The FFT does not reduce the number of partial amplitudes, and it does not
usually diagonalize the matrices inside each representation block.  It
replaces one arbitrary-looking $P$-by-$P$ problem by a collection of
structured block problems.  The result is exact; no colour approximation is
made.

It is best to measure the scaling in terms of $P$, the number of colour
orders, rather than writing $n!$ loosely.  For $n$ total gluons,
\begin{center}
\begin{tabular}{lll}
  \toprule
  implementation & permutation group & number of orders $P$ \\
  \midrule
  trace & $S_{n-1}$ & $(n-1)!$ \\
  adjoint/DDM & $S_{n-2}$ & $(n-2)!$ \\
  \bottomrule
\end{tabular}
\end{center}
Thus the adjoint calculation starts with exactly a factor $n-1$ fewer orders
than the trace calculation at the same multiplicity.

\begin{takeaway}
The dense $P^2$ colour matrix disappears.  What remains is a length-$P$
kernel, a fast transform of a length-$P$ amplitude vector, and contractions
inside Young-representation blocks.
\end{takeaway}

\section{Where the colour matrix comes from}

\subsection{Colour factors and partial amplitudes}

At fixed external momenta and helicities, a full gluon amplitude can be
written schematically as
\begin{equation}
  \cM = \sum_{g} \cC_g\,\cA(g).
  \label{eq:colour-decomposition}
\end{equation}
Here:
\begin{itemize}
  \item $g$ labels an ordering of some external gluons;
  \item $\cC_g$ is a colour tensor for that ordering;
  \item $\cA(g)$ is the corresponding complex partial amplitude.
\end{itemize}
The label $g$ will be treated as an element of a permutation group.  There
are $P$ possible values of $g$.

To obtain the full-colour squared amplitude, the external colour indices are
summed.  Inserting Eq.~\eqref{eq:colour-decomposition} gives
\begin{align}
  Q
  &= \sum_{\text{colours}} |\cM|^2 \\
  &= \sum_{g,h} \cA(g)^*\,
     \underbrace{\left\langle \cC_g,\cC_h\right\rangle}_{G_{g h}}
     \cA(h).
  \label{eq:gram-sum}
\end{align}
The matrix
\begin{equation}
  G_{g h}=\left\langle\cC_g,\cC_h\right\rangle
\end{equation}
is the colour Gram matrix.  It measures how much two colour tensors overlap.

In an orthonormal basis, $G$ would be the identity and the answer would be
just a sum of absolute squares.  That is what happens in the normalized
multiplet implementation.  Trace and DDM colour tensors are not orthogonal,
so their off-diagonal overlaps matter.

\subsection{The naive bottleneck}

For $P$ orderings, a generic Gram matrix contains $P^2$ numbers.  A direct
evaluation of Eq.~\eqref{eq:gram-sum} also performs order $P^2$ pairwise
interactions.  This becomes severe because $P$ is already factorial in the
number of gluons.

For example, the trace code has $P=(n-1)!$.  At eight total gluons this is
$7!=5040$ orders, so a dense matrix would have more than $25$ million
entries.  At ten gluons it would have
$(9!)^2$, more than $10^{11}$ entries.  The problem is not merely a
constant-factor optimization.

\section{The hidden structure: only the relative permutation matters}

The Gram matrix is not generic.  Relabelling both colour tensors by the same
permutation cannot change their fully contracted overlap.  In group notation,
\begin{equation}
  G_{qg,qh}=G_{g,h}
  \qquad\text{for every permutation }q.
  \label{eq:left-invariance}
\end{equation}
Choose $q=g^{-1}$.  Then
\begin{equation}
  G_{g,h}=G_{e,g^{-1}h}\equiv k(g^{-1}h),
  \label{eq:relative-kernel}
\end{equation}
where $e$ is the identity permutation and $k$ is a function on the group.

This statement has a simple interpretation.  The overlap does not care about
the absolute names of the two seating plans; it cares only about the
reshuffling required to turn one seating plan into the other.  Every row of
the large Gram matrix is therefore a rearranged copy of the same list of
kernel values.

Equation~\eqref{eq:relative-kernel} immediately reduces the colour data from
$P^2$ entries to $P$ entries.  It does not yet reduce a naive application
of the matrix.  Define
\begin{equation}
  B(h)=\sum_g \cA(g)\,k(g^{-1}h).
  \label{eq:convolution}
\end{equation}
Then
\begin{equation}
  Q=\sum_h \cA(h)^* B(h).
  \label{eq:q-from-convolution}
\end{equation}
Computing Eq.~\eqref{eq:convolution} directly is still a double sum over
$P^2$ pairs.  The purpose of the FFT is to evaluate this structured
convolution without treating those pairs independently.

\begin{takeaway}
The symmetry of the Gram matrix is the reason a fast method can exist.  The
FFT exploits that symmetry; it does not create it.
\end{takeaway}

\section{Analogy with the ordinary FFT}

Consider first a list indexed by points on a circle.  A circulant matrix has
entries that depend only on the relative displacement between two points.
Ordinary Fourier modes turn a circulant matrix into a diagonal matrix.  A
convolution then becomes multiplication frequency by frequency.

The present problem has the same architecture, with one difference:
\begin{itemize}
  \item positions on a circle form a cyclic group;
  \item gluon reorderings form a symmetric group $S_m$.
\end{itemize}
The cyclic group is abelian: doing two shifts in either order gives the same
result.  A symmetric group is non-abelian.  For example, if
$u=(12)$ and $v=(23)$, then
\begin{equation}
  uv\ne vu.
\end{equation}
Because permutation operations do not commute, their Fourier ``frequencies''
cannot all be represented by single complex numbers.  Some frequencies must
be represented by matrices.  Consequently, the permutation FFT generally
produces a \emph{block-diagonal} problem rather than a completely diagonal
one.

The analogy can be summarized as follows:
\begin{center}
\begin{tabular}{lll}
  \toprule
   & ordinary FFT & symmetric-group FFT \\
  \midrule
  index set & points on a cycle & permutations \\
  symmetry group & cyclic, abelian & $S_m$, non-abelian \\
  Fourier sector & one complex number & a small matrix \\
  structured matrix & circulant & group convolution \\
  transformed form & diagonal & block diagonal \\
  \bottomrule
\end{tabular}
\end{center}

No physical time or frequency is involved.  ``Frequency'' is only an analogy
for a component with a definite transformation pattern under permutations.

\section{The minimum representation theory we need}

\subsection{Representations and irreducible blocks}

A matrix representation assigns a matrix $\rho(g)$ to every group element
$g$, while preserving multiplication:
\begin{equation}
  \rho(gh)=\rho(g)\rho(h).
  \label{eq:representation}
\end{equation}
A representation is \emph{irreducible} if no change of basis can split all
its matrices into smaller independent blocks.  Irreducible representations
are the elementary symmetry sectors of the group.

For the symmetric group $S_m$, irreducible representations are labelled by
partitions of $m$, usually drawn as Young diagrams.  We denote a partition
by $\lambda$ and the dimension of its representation by $d_\lambda$.
The dimensions obey the crucial counting identity
\begin{equation}
  \sum_{\lambda\vdash m} d_\lambda^2=m!=|S_m|.
  \label{eq:sum-of-squares}
\end{equation}
Thus a function containing $m!$ scalar values can be reorganized into one
$d_\lambda$-by-$d_\lambda$ matrix for every Young diagram without losing
or adding information.

For $S_3$, the partitions and dimensions are
\begin{center}
\begin{tabular}{lll}
  \toprule
  partition & informal symmetry & dimension \\
  \midrule
  $(3)$ & completely symmetric & $1$ \\
  $(2,1)$ & mixed symmetry & $2$ \\
  $(1,1,1)$ & completely antisymmetric & $1$ \\
  \bottomrule
\end{tabular}
\end{center}
Indeed,
\begin{equation}
  1^2+2^2+1^2=6=|S_3|.
\end{equation}

\subsection{The transform used by the code}

The implementation uses the unnormalized forward transform
\begin{equation}
  \widehat f_\lambda
  =\sum_{g\in S_m}f(g)\rho_\lambda(g).
  \label{eq:forward-transform}
\end{equation}
For a one-dimensional representation, $\widehat f_\lambda$ is one complex
number.  For a representation of dimension $d_\lambda$, it is a
$d_\lambda$-by-$d_\lambda$ complex matrix.

There are still exactly $P=m!$ scalar Fourier coefficients in total, by
Eq.~\eqref{eq:sum-of-squares}.  The FFT is therefore not a lossy compression.
It is a reorganization that exposes which pieces can and cannot interact.

\section{Why convolution becomes block multiplication}

Let $B$ be the colour convolution in Eq.~\eqref{eq:convolution}.  Set
$x=g^{-1}h$, so that $h=gx$.  Transforming $B$ gives
\begin{align}
  \widehat B_\lambda
  &=\sum_h B(h)\rho_\lambda(h) \\
  &=\sum_{g,x}\cA(g)k(x)\rho_\lambda(gx) \\
  &=\left(\sum_g\cA(g)\rho_\lambda(g)\right)
    \left(\sum_x k(x)\rho_\lambda(x)\right) \\
  &=\widehat A_\lambda K_\lambda,
  \label{eq:convolution-theorem}
\end{align}
where
\begin{equation}
  \widehat A_\lambda=\sum_g\cA(g)\rho_\lambda(g),
  \qquad
  K_\lambda=\sum_x k(x)\rho_\lambda(x).
\end{equation}

This is the non-abelian convolution theorem.  A sum coupling every pair of
permutations has become one ordinary matrix product in each Young sector.
Different Young diagrams never mix.

The finite-group version of Parseval's or Plancherel's identity then converts
Eq.~\eqref{eq:q-from-convolution} into
\begin{equation}
  Q=\frac{1}{|S_m|}
    \sum_{\lambda\vdash m}d_\lambda
    \Tr\!\left(
      \widehat A_\lambda^\dagger
      \widehat A_\lambda K_\lambda
    \right).
  \label{eq:block-quadratic-form}
\end{equation}
This is the central equation implemented by both colour codes.

The factors in Eq.~\eqref{eq:block-quadratic-form} have straightforward
roles:
\begin{itemize}
  \item $1/|S_m|$ compensates for the unnormalized forward transform;
  \item $d_\lambda$ is the Plancherel weight of the irrep;
  \item $\widehat A_\lambda^\dagger\widehat A_\lambda$ contains overlaps
        between columns of the transformed amplitude block;
  \item $K_\lambda$ supplies the colour weight for those overlaps.
\end{itemize}

For the colour kernels in these codes, inversion symmetry of $k$ makes each
$K_\lambda$ real symmetric, up to floating-point roundoff.  The
implementation checks this property and stores the symmetrized result.  Its
contraction evaluates
\begin{equation}
  \Tr(\widehat A^\dagger\widehat A K)
  =\sum_i K_{ii}\|\widehat A_i\|^2
   +2\sum_{i<j}K_{ij}
      \operatorname{Re}(\widehat A_i^\dagger\widehat A_j),
  \label{eq:symmetric-block-contraction}
\end{equation}
where $\widehat A_i$ is column $i$ of an amplitude block.  Each
off-diagonal column overlap is therefore computed only once.

\begin{takeaway}
The FFT has not thrown away the operation $\cA^\dagger G\cA$.  It evaluates
the same scalar in a basis where unrelated symmetry sectors are already
separated.
\end{takeaway}

\section{Warm-up example: the four-gluon DDM matrix}

At four gluons, the DDM basis contains two orders.  In the normalization used
by \textsc{AmpliGluonAdjoint}, its Gram matrix is
\begin{equation}
  G=\begin{pmatrix}
       288 & 144\\
       144 & 288
     \end{pmatrix}.
  \label{eq:four-gluon-ddm-gram}
\end{equation}
Let the two partial amplitudes be $a$ and $b$.  Direct multiplication gives
\begin{equation}
  Q=288\bigl(|a|^2+|b|^2\bigr)
    +288\operatorname{Re}(a^*b).
  \label{eq:four-gluon-ddm-direct}
\end{equation}

The two orders form $S_2$.  Its two one-dimensional Fourier sectors are the
symmetric and antisymmetric combinations
\begin{equation}
  \widehat A_+=a+b,
  \qquad
  \widehat A_-=a-b.
\end{equation}
Writing $\tau=(12)$ for the nontrivial permutation, the kernel values
$k(e)=288$ and $k(\tau)=144$ transform to
\begin{equation}
  K_+=432,
  \qquad
  K_-=144.
\end{equation}
Equation~\eqref{eq:block-quadratic-form} becomes
\begin{equation}
  Q=\frac{1}{2}
    \left(432|a+b|^2+144|a-b|^2\right),
  \label{eq:four-gluon-ddm-fourier}
\end{equation}
which expands to Eq.~\eqref{eq:four-gluon-ddm-direct}.

This tiny example is fully diagonal because $S_2$ is abelian and all its
irreps are one-dimensional.  For larger symmetric groups, mixed-symmetry
irreps have dimension greater than one and produce matrix blocks.

\section{Worked block example: the four-gluon trace kernel on
\texorpdfstring{$S_3$}{S3}}

The four-gluon trace code fixes one gluon and permutes the other three, giving
six orders in $S_3$.  In the code's $U(3)$ loop normalization, the kernel
is
\begin{equation}
  k(e)=81,
  \qquad
  k(g)=9\quad\text{for }g\ne e.
  \label{eq:s3-trace-kernel}
\end{equation}

Write
\begin{equation}
  S_3=\{e,r,r^2,s,sr,sr^2\},
  \qquad r^3=s^2=e,
  \qquad srs=r^{-1}.
\end{equation}
Besides the trivial and sign representations, $S_3$ has a standard
two-dimensional representation.  One real choice is
\begin{equation}
  \rho_{\mathrm{std}}(r)=
  \begin{pmatrix}
    -\tfrac12&-\tfrac{\sqrt3}{2}\\
     \tfrac{\sqrt3}{2}&-\tfrac12
  \end{pmatrix},
  \qquad
  \rho_{\mathrm{std}}(s)=
  \begin{pmatrix}1&0\\0&-1\end{pmatrix}.
\end{equation}

Transforming Eq.~\eqref{eq:s3-trace-kernel} gives three blocks:
\begin{equation}
  K_{\mathrm{triv}}=126,
  \qquad
  K_{\mathrm{sign}}=72,
  \qquad
  K_{\mathrm{std}}=72I_2.
  \label{eq:s3-kernel-blocks}
\end{equation}
The information count is $1+1+4=6$, exactly the number of original kernel
values.

For a transparent numerical check, suppose only two amplitude values are
nonzero:
\begin{equation}
  \cA(e)=a,
  \qquad
  \cA(s)=b.
\end{equation}
Their Fourier blocks are
\begin{align}
  \widehat A_{\mathrm{triv}}&=a+b,\\
  \widehat A_{\mathrm{sign}}&=a-b,\\
  \widehat A_{\mathrm{std}}
    &=aI_2+b\rho_{\mathrm{std}}(s)
      =\begin{pmatrix}a+b&0\\0&a-b\end{pmatrix}.
\end{align}
The block formula gives
\begin{align}
  Q=\frac{1}{6}\Bigl[&126|a+b|^2+72|a-b|^2 \\
      &+2(72)\bigl(|a+b|^2+|a-b|^2\bigr)\Bigr] \\
   &=81\bigl(|a|^2+|b|^2\bigr)
     +18\operatorname{Re}(a^*b).
  \label{eq:s3-block-answer}
\end{align}
The direct two-entry calculation uses $G_{ee}=G_{ss}=81$ and
$G_{es}=G_{se}=9$, and gives the same result.  For $a=b=1$, both methods
give $Q=180$.

This four-gluon kernel is unusually simple: its two-dimensional block happens
to be proportional to the identity.  At higher multiplicity, the code allows
and contracts full real symmetric matrices inside the Young blocks.  The
general claim is block diagonalization, not complete diagonalization.

\section{Why this is a \emph{fast} group Fourier transform}

Equation~\eqref{eq:forward-transform} defines a Fourier transform, but using
that definition literally would not be fast.  One would construct many dense
matrices $\rho_\lambda(g)$ and sum over every permutation separately.  The
implementation instead reuses work along the subgroup tower
\begin{equation}
  S_1\subset S_2\subset\cdots\subset S_m.
  \label{eq:subgroup-tower}
\end{equation}

\subsection{Recursive permutation order}

Let $s_i=(i,i+1)$ be an adjacent swap.  The code organizes $S_m$ into
$m$ cosets of $S_{m-1}$:
\begin{equation}
  S_m=\bigsqcup_{i=1}^{m}
      (s_i s_{i+1}\cdots s_{m-1})S_{m-1},
  \label{eq:coset-decomposition}
\end{equation}
where the product is empty when $i=m$.  In elementary terms, first choose
which label occupies the final image position, and then permute the remaining
$m-1$ labels.  This produces $m$ copies of the same $S_{m-1}$ problem.

The input amplitudes arrive in lexicographic permutation order.  A one-time
index map rearranges them into the recursive order of
Eq.~\eqref{eq:coset-decomposition}.  The transform can then reuse the already
computed $S_{m-1}$ transforms for every coset.

\subsection{Young branching}

When an irrep of $S_m$ is restricted to $S_{m-1}$, it splits according to
the removable corners of its Young diagram.  In the standard Young-tableau
basis, these child representations occupy diagonal subblocks.  The code can
therefore copy each transformed $S_{m-1}$ child directly into the
appropriate part of its $S_m$ parent block.

No representation-theory expertise is needed to understand the computational
point: the chosen basis makes the recursive subproblem visibly block
diagonal, so results from the previous level can be reused rather than
recomputed.

\subsection{Sparse adjacent swaps}

The remaining coset representative is a product of adjacent swaps.  In
Young's orthogonal representation, an adjacent swap either
\begin{itemize}
  \item multiplies one tableau state by a sign, or
  \item mixes a pair of tableau states through a $2$-by-$2$ orthogonal block.
\end{itemize}
If the axial distance for the first tableau is $a$, the two-state block can
be written as
\begin{equation}
  \begin{pmatrix}
    \dfrac{1}{a} & \sqrt{1-\dfrac{1}{a^2}} \\
    \sqrt{1-\dfrac{1}{a^2}} & -\dfrac{1}{a}
  \end{pmatrix}.
\end{equation}
This matrix is a reflection: it is symmetric and its square is the identity.
The implementation stores, for each state and generator, only its partner,
the diagonal coefficient, and the mixing coefficient.  It never needs a
general dense generator matrix.

Combining recursive reuse, Young branching, and sparse adjacent swaps is what
makes the transform an FFT rather than a direct group Fourier transform.

\section{How \textsc{AmpliGluonTrace} uses the FFT}

\subsection{The ordering space}

For $n$ total gluons, the trace code fixes the final gluon and retains all
orders of the first $n-1$.  Its amplitude vector is therefore a function on
$S_{n-1}$, with
\begin{equation}
  P_{\mathrm{trace}}=(n-1)!.
\end{equation}
These trace tensors are non-orthogonal, so a colour Gram contraction is
required.

\subsection{The trace kernel}

The kernel is most easily constructed using fundamental colour lines.  Embed
$h\in S_{n-1}$ into $S_n$ by making it fix $n$, define the long cycle
$\gamma=(1\,2\,\ldots\,n)$, and let $\nu(\pi)$ count the cycles of a
permutation $\pi$.  The implemented $U(3)$ kernel is
\begin{equation}
  k_{\mathrm{trace}}(h)
  =3^{\nu(\gamma^{-1}h\gamma h^{-1})}.
  \label{eq:trace-loop-kernel}
\end{equation}
Every closed fundamental-index loop supplies one factor of three.  For a
pure-gluon tree amplitude, the $U(1)$ component decouples, so this $U(3)$
construction yields the same contracted result as the required $SU(3)$
colour sum.

\subsection{One-time and per-event work}

For a fixed multiplicity, the colour object performs the following setup:
\begin{enumerate}
  \item build the $P_{\mathrm{trace}}$ kernel values;
  \item initialize the $S_{n-1}$ Fourier plan and reusable workspace;
  \item transform the kernel into the blocks $K_\lambda$;
  \item verify that every transformed kernel block is real symmetric.
\end{enumerate}
For each nonzero amplitude evaluation, it then:
\begin{enumerate}
  \item fills all ordered partial amplitudes $\cA(g)$;
  \item transforms them into $\widehat A_\lambda$;
  \item evaluates Eq.~\eqref{eq:block-quadratic-form} block by block.
\end{enumerate}
The transformed kernel is reused.  Only the amplitudes depend on the current
phase-space point and helicity assignment.

\section{How \textsc{AmpliGluonAdjoint} uses the FFT}

\subsection{The smaller DDM ordering space}

The DDM decomposition fixes two gluons at the ends of a chain and permutes
the $n-2$ gluons between them:
\begin{equation}
  \cM=\sum_{\sigma\in S_{n-2}}
      \cC_\sigma\,\cA(1,\sigma,n).
\end{equation}
Consequently,
\begin{equation}
  P_{\mathrm{DDM}}=(n-2)!.
\end{equation}
The DDM basis is smaller than the full trace ordering set, but it is still
non-orthogonal.  Its Gram entries again depend only on the relative DDM
permutation, so the colour matrix is a convolution on $S_{n-2}$.

\subsection{Building the adjoint kernel without expanding every chain}

A DDM colour factor is a chain of adjoint structure constants.  In the
fundamental representation, that chain is a nested commutator.  Every
commutator
\begin{equation}
  [X,Y]=XY-YX
\end{equation}
is a finite difference between two trace orderings.  Expanding all nested
commutators in both colour factors independently would generate a large
number of trace terms.

The adjoint code instead starts with the trace kernel and applies the
commutators as repeated finite-difference sweeps over the permutation array.
The left and right DDM chains require $2(n-2)$ such sweeps over the
$S_{n-1}$ trace kernel.  The final restricted array contains the
$(n-2)!$ values of the DDM kernel.  This factorization avoids separately
forming all trace pairings for every DDM Gram entry.

The use of the $U(3)$ trace kernel is exact here as well: the identity
generator commutes with everything, so every $U(1)$ contribution vanishes
inside the commutator chain.

\subsection{The transform and contraction are shared}

After constructing its kernel, the adjoint code uses the same
\code{symmetric\_group\_fft} implementation and the same block formula as the
trace code, now at degree $n-2$.  In particular:
\begin{itemize}
  \item the kernel is transformed once and stored as real symmetric blocks;
  \item each DDM amplitude vector is transformed on evaluation;
  \item the result is accumulated using Eq.~\eqref{eq:symmetric-block-contraction}.
\end{itemize}
The two codes therefore differ in their amplitude bases, group degrees, and
kernel construction, but not in the Fourier mathematics used after the
kernel is available.

\section{What scaling is actually improved?}

This is where precise language matters.  Let $P=|S_m|=m!$ denote the number
of orders in whichever code is being discussed.

\subsection{Storage}

The storage improvement is clean:
\begin{center}
\begin{tabular}{ll}
  \toprule
  object & number of stored scalars \\
  \midrule
  dense Gram matrix & $P^2$ \\
  relative-permutation kernel & $P$ \\
  all Fourier kernel blocks &
     $\sum_\lambda d_\lambda^2=P$ \\
  transformed amplitude blocks &
     $\sum_\lambda d_\lambda^2=P$ \\
  \bottomrule
\end{tabular}
\end{center}
Thus the colour-matrix storage is reduced from quadratic to linear in the
number of orders.

\subsection{Arithmetic}

A direct dense quadratic form takes order $P^2$ pair operations.  The
Fourier route replaces it by
\begin{enumerate}
  \item one fast group transform of the amplitude vector;
  \item one contraction for each Young block.
\end{enumerate}
The recursive FFT is designed to cost roughly a factorial-sized traversal
times multiplicity-dependent polynomial overhead, rather than a direct
all-permutations-by-all-Fourier-coefficients calculation.

For a block of dimension $d_\lambda$, the implemented column-overlap
contraction takes order $d_\lambda^3$ arithmetic operations.  The total
block work is therefore characterized by
\begin{equation}
  \sum_\lambda d_\lambda^3.
  \label{eq:block-cost}
\end{equation}
For every nontrivial case, this is bounded well below the $P^2$ pairwise cost
of the original dense quadratic form.  A simple, deliberately loose bound is
\begin{equation}
  \sum_\lambda d_\lambda^3
  \leq
  \left(\max_\lambda d_\lambda\right)
  \sum_\lambda d_\lambda^2
  \leq P^{3/2}.
  \label{eq:block-bound}
\end{equation}
The real-symmetric kernel lets the implementation evaluate each off-diagonal
column overlap once instead of twice.  The implementation also tiles several
overlaps together, improving constants but not changing this characterization.

Consequently, the safest statement is:
\begin{quote}
The algorithm removes the $P^2$ dense colour-matrix bottleneck.  It replaces
it with $P$-sized colour data, a fast permutation transform, and independent
subquadratic block contractions.
\end{quote}
Calling this ``about $P$,'' and hence ``about factorial,'' is a useful
high-level summary of the storage and transform structure.  It is not a proof
that the complete runtime is exactly $\Theta(P)$.

\subsection{The whole amplitude remains factorial}

Neither FFT reduces the number of partial amplitudes:
\begin{equation}
  P_{\mathrm{trace}}=(n-1)!,
  \qquad
  P_{\mathrm{DDM}}=(n-2)!.
\end{equation}
Those amplitudes must still be generated, although the codes share
Berends--Giele subcurrents and use special MHV evaluators to reduce their
kinematic cost.  The colour FFT addresses the colour contraction stage, not
the entire factorial growth of an ordered-amplitude formulation.

Kernel construction is also separate from per-event contraction.  It is done
once per multiplicity and can be amortized over many phase-space points and
helicity configurations.

\begin{takeaway}
``Factorial squared to about factorial'' correctly identifies the bottleneck
that disappears, provided it is not read as a strict asymptotic theorem for
the entire program.
\end{takeaway}

\section{How established is this method?}

The answer depends on which part of the method one has in mind.  It is useful
to distinguish three layers:
\begin{enumerate}
  \item The theorem that a group convolution becomes block multiplication
        after a finite-group Fourier transform is classical.
  \item Fast algorithms specialized to the symmetric group are an established
        subject in computational representation theory, but they are much less
        widely known than the ordinary FFT.
  \item Applying such a fast numerical transform to a trace- or DDM-basis
        colour Gram contraction appears to be uncommon in collider QCD.
\end{enumerate}
Thus the mathematical identity is not new.  What is unusual here is the way
in which it is used to remove a concrete collider-calculation bottleneck.

\subsection{Classical mathematics versus an actual FFT}

For a finite group, a matrix of the form
\begin{equation}
  G_{\sigma,\tau}=k(\sigma^{-1}\tau)
\end{equation}
is a group-convolution operator.  Decomposing the group algebra into its
irreducible matrix blocks is the finite-group analogue of diagonalizing a
circulant matrix with an ordinary Fourier transform.  For a non-abelian group
such as $S_m$, the result is generally block diagonal rather than scalar
diagonal.  Only a central kernel, also called a class function, is guaranteed
to act as a scalar within every irreducible block.

There is an important algorithmic distinction.  Equation
\eqref{eq:forward-transform} defines Fourier coefficients by summing over all
$m!$ permutations.  Evaluating that definition independently for every
coefficient is still an order-$(m!)^2$ calculation.  A genuine FFT exploits
the subgroup tower in Eq.~\eqref{eq:subgroup-tower}, Young branching, and
sparse adjacent-transposition matrices to reuse intermediate results.

Fast transforms of this kind were developed by Clausen and by Diaconis and
Rockmore, with subsequent refinements and implementations by Clausen and Baum
and by Maslen~\cite{clausen1989,diaconisrockmore1990,clausenbaum1993,maslen1998}.
Maslen gives an explicit bound of at most
\begin{equation}
  \frac{3}{4}m(m-1)|S_m|
\end{equation}
multiplications, and the same number of additions, for his symmetric-group
transform.  The appropriate summary is therefore ``linear in the group order
up to a polynomial in $m$,'' rather than literally $O(m!)$ with a
multiplicity-independent constant.  This is the established algorithmic
precedent for the FFT used in this repository.

Merely labelling the rows and columns of an arbitrary matrix by permutations
does not create this speedup.  The relative-permutation identity above, or a
closely related subgroup or coset symmetry, is the essential hypothesis.

\subsection{Use in formal high-energy theory}

Symmetric-group Fourier ideas are well established in formal high-energy
theory.  Multi-trace operators containing many matrix fields are naturally
described by permutations.  Passing from a permutation basis to a basis
labelled by symmetric-group representations leads to Schur and restricted
Schur operators, organizes finite-$N_c$ relations, and often diagonalizes
free two-point functions.  Early and representative examples include the
work of Corley, Jevicki, and Ramgoolam and of Brown, Heslop, and
Ramgoolam~\cite{Corley:2001zk,Brown:2007xh}.  Broader treatments describe the
Fourier transformation of permutation algebras in gauge-invariant operator
bases, quiver gauge theories, and tensor models
~\cite{Ramgoolam:2016ciq,Mattioli:2016gyl,BenGeloun:2017vwn}.

These applications use the same underlying representation-theory
decomposition, but usually for a different purpose.  They use the Fourier
basis analytically to classify operators or diagonalize correlators.  They do
not generally claim to run a recursive Clausen--Maslen FFT in order to speed
up repeated numerical convolutions.  The words ``Fourier transform'' in a
formal field-theory paper therefore do not by themselves imply the fast
algorithm described by Eq.~\eqref{eq:subgroup-tower}.

\subsection{Status in collider QCD}

Collider calculations certainly use permutation-labelled trace and DDM
bases, but their common computational strategies are different.  Examples
include colour-flow decompositions, colour-dressed recursion, and orthogonal
$SU(N_c)$ multiplet bases~\cite{Maltoni:2002mq,Duhr:2006iq,Keppeler:2012ih}.  The
\textsc{ColorFull} library is an example of direct trace-basis colour
algebra~\cite{Sjodahl:2014opa}.

To our knowledge, a fast numerical $S_m$ transform has not become a standard
method for collider colour sums.  In particular, we are not aware of an
earlier collider implementation that identifies the trace or DDM Gram matrix
as a relative-permutation convolution and then uses a recursive
symmetric-group FFT to avoid the dense $\cA^\dagger G\cA$ contraction.  This
is deliberately a cautious statement about the literature, not a categorical
claim of priority.

The most accurate characterization is therefore:
\begin{quote}
The finite-group convolution theorem and fast symmetric-group transforms are
classical.  Their use as analytic representation bases is established in
formal gauge theory and tensor models.  Their use as a numerical accelerator
for collider-QCD colour Gram contractions is apparently uncommon.
\end{quote}

\subsection{Where symmetric-group Fourier analysis is used most}

The most developed practical literature concerns data whose states are
themselves permutations.  Diaconis used symmetric-group harmonic analysis
for ranked data~\cite{diaconis1989}.  Later work applied fast transforms to
probability distributions over assignments, multi-object tracking, and
clustering of rankings~\cite{kondor2007,plumb2015}.  There is also a quantum
algorithm for the Fourier transform over symmetric groups~\cite{beals1997}.
These fields treat the transform as a specialized tool rather than a
universal primitive comparable in popularity to the ordinary FFT.

A closely related physics application is the Fourier analysis of
many-particle transition amplitudes for partially distinguishable bosons and
fermions~\cite{Dufour:2024nmz}.  There too, permutations label the different ways
particles can be assigned between input and output states.  This resembles
the ordering structure of the colour problem, but it is not a QCD colour
calculation.

\subsection{Does the application require all-gluon scattering?}

No general theorem requires gluons.  The transform applies whenever the
large matrix really is a group convolution.  All-gluon trace and DDM bases
are especially clean because the relevant orderings furnish a full symmetric
group, and simultaneous relabelling makes the Gram entry depend only on the
relative permutation.

With quarks or several distinguishable particle species, fixed endpoints and
species labels can reduce the symmetry.  The remaining structure may be a
convolution on a subgroup, a function on a coset space, or a collection of
smaller convolution problems.  Fourier methods can still apply, but the
present $S_m$ implementation may need to be adapted.  Conversely, an
arbitrary matrix whose indices merely happen to be permutations receives no
benefit unless its entries possess the required relative-permutation
invariance.

\section{Comparison with the multiplet basis}

The three colour strategies in this repository solve related but different
linear-algebra problems:
\begin{center}
\small
\begin{tabular}{@{}
  >{\raggedright\arraybackslash}p{0.18\linewidth}
  >{\raggedright\arraybackslash}p{0.21\linewidth}
  >{\raggedright\arraybackslash}p{0.22\linewidth}
  >{\raggedright\arraybackslash}p{0.30\linewidth}@{}}
  \toprule
  implementation & basis & Gram matrix & final colour operation \\
  \midrule
  multiplet & normalized multiplet paths & identity & Euclidean norm \\
  trace & single traces & non-diagonal convolution & FFT plus Young-block contraction \\
  adjoint & DDM chains & non-diagonal convolution & FFT plus Young-block contraction \\
  \bottomrule
\end{tabular}
\end{center}

The permutation FFT does not secretly turn the trace or DDM tensors into the
multiplet basis.  It decomposes the \emph{space of ordering functions} into
irreducible sectors of a symmetric group.  Within a sector, colour components
can still mix.  The multiplet implementation instead starts from colour
tensors whose Gram matrix is already the identity.

\section{Algorithm summary}

The common logic can be written as the following pseudocode.

\subsection*{Setup for one gluon multiplicity}

\begin{enumerate}
  \item Choose $m=n-1$ for trace or $m=n-2$ for DDM.
  \item Build the $P=m!$ relative-permutation kernel values $k(g)$.
  \item Enumerate Young diagrams and standard tableaux for every
        $S_1,\ldots,S_m$.
  \item Build sparse adjacent-swap representations and the recursive coset
        ordering.
  \item Compute $K_\lambda=\operatorname{FFT}_{S_m}[k]_\lambda$.
  \item Check that every $K_\lambda$ is finite, real, and symmetric.
\end{enumerate}

\subsection*{For each nonzero amplitude vector}

\begin{enumerate}
  \item Compute all required partial amplitudes $\cA(g)$.
  \item Compute
        $\widehat A_\lambda=\operatorname{FFT}_{S_m}[\cA]_\lambda$.
  \item For every Young diagram, accumulate
        $d_\lambda\Tr(
          \widehat A_\lambda^\dagger
          \widehat A_\lambda K_\lambda)$.
  \item Divide the sum by $m!$.
  \item Apply any separately requested initial-colour average.
\end{enumerate}

Exact-zero helicity sectors and zero coupling can return before the transform,
which is an independent optimization.

\section{Common misconceptions}

\paragraph{``The FFT makes the original basis orthogonal.''}
No.  Trace and DDM colour tensors remain non-orthogonal.  The transform makes
their convolution structure block diagonal in permutation-representation
space.

\paragraph{``Every Fourier block is diagonal.''}
No.  That is true for abelian groups and happens accidentally in some small
examples.  General $S_m$ irreps are multidimensional, and their kernel
blocks are matrices.

\paragraph{``The FFT reduces the number of partial amplitudes.''}
No.  It transforms all $P$ amplitudes into exactly $P$ scalar Fourier
coefficients arranged in blocks.

\paragraph{``This is an approximation to full colour.''}
No.  For the physical pure-gluon amplitudes in the scope of these codes, the
convolution theorem, the $U(1)$ decoupling used by the trace kernel, and the
Plancherel identity are exact.  Differences from a direct dense contraction
are only floating-point roundoff.

\paragraph{``Young diagrams describe gluon species.''}
No.  They label transformation patterns under permutations of ordering
positions.  They are bookkeeping devices for permutation symmetry.

\paragraph{``The adjoint commutators and the non-abelian FFT are the same
noncommutativity.''}
No.  The commutators construct DDM colour factors in the fundamental colour
algebra.  The matrix-valued FFT arises because permutations themselves do not
commute.  These are separate structures that happen to meet in the same
calculation.

\paragraph{``The runtime is now proven to be exactly factorial.''}
No.  The stored colour data are linear in $P$, and the dense $P^2$
bottleneck is gone.  Transform overhead, Young-block dimensions, and the cost
of producing the partial amplitudes remain.

\section{Map from the explanation to the source code}

The relevant implementation is compactly divided by responsibility:
\begin{center}
\small
\begin{tabular}{@{}
  >{\raggedright\arraybackslash}p{0.43\linewidth}
  >{\raggedright\arraybackslash}p{0.52\linewidth}@{}}
  \toprule
  source file & role \\
  \midrule
  \path{AmpliGluonTrace/src/symmetric_group_fft.f90}
    & Young diagrams and tableaux, adjacent generators, subgroup recursion,
      forward transform, and reusable workspace \\
  \path{AmpliGluonTrace/src/trace_colour_kernel.f90}
    & cycle-count construction of the $U(3)$ trace kernel \\
  \path{AmpliGluonTrace/src/trace_colour_matrix.f90}
    & kernel transform, amplitude transform, and blockwise Plancherel
      contraction for trace orders \\
  \path{AmpliGluonAdjoint/src/adjoint_colour_matrix.f90}
    & finite-difference construction of the DDM kernel and its blockwise
      contraction \\
  \path{AmpliGluonAdjoint/Makefile}
    & compiles the same shared FFT and trace-kernel sources into the adjoint
      library \\
  \bottomrule
\end{tabular}
\end{center}

Useful regression tests are found in
\path{AmpliGluonTrace/tests/test_symmetric_group_fft.f90} and
\path{AmpliGluonTrace/tests/test_trace_colour_matrix.f90}.  They compare the
fast transform with direct representation sums, check Parseval identities,
and compare the Fourier colour contraction with direct dense convolution on
small groups.

\section{Conclusion}

The essential chain of reasoning is now short:
\begin{enumerate}
  \item A trace or DDM colour decomposition produces a non-diagonal Gram
        matrix.
  \item Simultaneous relabelling invariance makes each Gram entry depend only
        on a relative permutation.
  \item The Gram matrix is therefore a convolution on $S_{n-1}$ or
        $S_{n-2}$, specified by one factorial-sized kernel.
  \item Fourier transformation on the symmetric group separates the
        convolution into Young-irrep blocks.
  \item A recursive subgroup algorithm and sparse adjacent swaps make that
        transform fast enough to be useful.
  \item The transformed kernel is reused, while every event requires only an
        amplitude transform and independent block contractions.
\end{enumerate}

This is precisely how the codes avoid the explicit dense
$\cA^\dagger G\cA$ multiplication.  They do not change the physics and do
not discard interference terms.  They expose the permutation symmetry that
was already present, and use it to evaluate the same colour sum in a far more
structured representation.

\begin{thebibliography}{99}

\bibitem{clausen1989}
M.~Clausen,
``Fast generalized Fourier transforms,''
\emph{Theoretical Computer Science} \textbf{67} (1989) 55--63,
\href{https://doi.org/10.1016/0304-3975(89)90021-2}{doi:10.1016/0304-3975(89)90021-2}.

\bibitem{diaconisrockmore1990}
P.~Diaconis and D.~Rockmore,
``Efficient computation of the Fourier transform on finite groups,''
\emph{Journal of the American Mathematical Society} \textbf{3} (1990)
297--332,
\href{https://doi.org/10.1090/S0894-0347-1990-1030655-4}{doi:10.1090/S0894-0347-1990-1030655-4}.

\bibitem{clausenbaum1993}
M.~Clausen and U.~Baum,
``Fast Fourier transforms for symmetric groups: theory and implementation,''
\emph{Mathematics of Computation} \textbf{61} (1993) 833--847,
\href{https://doi.org/10.1090/S0025-5718-1993-1192969-X}{doi:10.1090/S0025-5718-1993-1192969-X}.

\bibitem{maslen1998}
D.~K.~Maslen,
``The efficient computation of Fourier transforms on the symmetric group,''
\emph{Mathematics of Computation} \textbf{67} (1998) 1121--1147,
\href{https://doi.org/10.1090/S0025-5718-98-00964-8}{doi:10.1090/S0025-5718-98-00964-8}.

\bibitem{Corley:2001zk}
S.~Corley, A.~Jevicki, and S.~Ramgoolam,
``Exact correlators of giant gravitons from dual $\mathcal N=4$ SYM theory,''
\emph{Advances in Theoretical and Mathematical Physics} \textbf{5} (2002)
809--839,
\href{https://arxiv.org/abs/hep-th/0111222}{arXiv:hep-th/0111222}.

\bibitem{Brown:2007xh}
T.~W.~Brown, P.~J.~Heslop, and S.~Ramgoolam,
``Diagonal multi-matrix correlators and BPS operators in $\mathcal N=4$ SYM,''
\emph{Journal of High Energy Physics} \textbf{02} (2008) 030,
\href{https://arxiv.org/abs/0711.0176}{arXiv:0711.0176}.

\bibitem{Ramgoolam:2016ciq}
S.~Ramgoolam,
``Permutations and the combinatorics of gauge invariants for general $N$,''
\href{https://arxiv.org/abs/1605.00843}{arXiv:1605.00843}.

\bibitem{Mattioli:2016gyl}
P.~Mattioli and S.~Ramgoolam,
``Gauge invariants and correlators in flavoured quiver gauge theories,''
\emph{Nuclear Physics B} \textbf{911} (2016) 638--711,
\href{https://arxiv.org/abs/1603.04369}{arXiv:1603.04369}.

\bibitem{BenGeloun:2017vwn}
J.~Ben Geloun and S.~Ramgoolam,
``Tensor models, Kronecker coefficients and permutation centralizer algebras,''
\emph{Journal of High Energy Physics} \textbf{11} (2017) 092,
\href{https://arxiv.org/abs/1708.03524}{arXiv:1708.03524}.

\bibitem{Maltoni:2002mq}
F.~Maltoni, K.~Paul, T.~Stelzer, and S.~Willenbrock,
``Color-flow decomposition of QCD amplitudes,''
\emph{Physical Review D} \textbf{67} (2003) 014026,
\href{https://arxiv.org/abs/hep-ph/0209271}{arXiv:hep-ph/0209271}.

\bibitem{Duhr:2006iq}
C.~Duhr, S.~Hoche, and F.~Maltoni,
``Color-dressed recursive relations for multi-parton amplitudes,''
\emph{Journal of High Energy Physics} \textbf{08} (2006) 062,
\href{https://arxiv.org/abs/hep-ph/0607057}{arXiv:hep-ph/0607057}.

\bibitem{Keppeler:2012ih}
S.~Keppeler and M.~Sjodahl,
``Orthogonal multiplet bases in $SU(N_c)$ color space,''
\emph{Journal of High Energy Physics} \textbf{09} (2012) 124,
\href{https://arxiv.org/abs/1207.0609}{arXiv:1207.0609}.

\bibitem{Sjodahl:2014opa}
M.~Sjodahl,
``ColorFull: a C++ library for calculations in $SU(N_c)$ color space,''
\emph{European Physical Journal C} \textbf{75} (2015) 236,
\href{https://arxiv.org/abs/1412.3967}{arXiv:1412.3967}.

\bibitem{diaconis1989}
P.~Diaconis,
``A generalization of spectral analysis with application to ranked data,''
\emph{Annals of Statistics} \textbf{17} (1989) 949--979,
\href{https://www.jstor.org/stable/2241705}{JSTOR 2241705}.

\bibitem{kondor2007}
R.~Kondor, A.~Howard, and T.~Jebara,
``Multi-object tracking with representations of the symmetric group,''
in \emph{Proceedings of AISTATS} (2007),
\href{https://proceedings.mlr.press/v2/kondor07a.html}{PMLR article}.

\bibitem{plumb2015}
G.~Plumb, B.~Pachauri, R.~Kondor, and V.~Singh,
``SnFFT: a Julia toolkit for Fourier analysis of functions over permutations,''
\emph{Journal of Machine Learning Research} \textbf{16} (2015) 3469--3473,
\href{https://jmlr.org/papers/v16/plumb15a.html}{JMLR article}.

\bibitem{beals1997}
R.~Beals,
``Quantum computation of Fourier transforms over symmetric groups,''
in \emph{Proceedings of the 29th Annual ACM Symposium on Theory of
Computing} (1997) 48--53,
\href{https://doi.org/10.1145/258533.258548}{doi:10.1145/258533.258548}.

\bibitem{Dufour:2024nmz}
G.~Dufour and A.~Buchleitner,
``Fourier analysis of many-body transition amplitudes and states,''
\href{https://arxiv.org/abs/2409.15079}{arXiv:2409.15079}.

\end{thebibliography}

\appendix

\section{A compact derivation of the Plancherel formula}

For unitary irreducible representations of a finite group $G$, matrix
elements satisfy the orthogonality relation
\begin{equation}
  \sum_{g\in G}
  \rho_\lambda(g)_{ij}^*
  \rho_\mu(g)_{k\ell}
  =\frac{|G|}{d_\lambda}
   \delta_{\lambda\mu}\delta_{ik}\delta_{j\ell}.
  \label{eq:matrix-element-orthogonality}
\end{equation}
This implies, for the unnormalized transform of
Eq.~\eqref{eq:forward-transform},
\begin{equation}
  \sum_{g\in G} f(g)^*q(g)
  =\frac{1}{|G|}
   \sum_\lambda d_\lambda
   \Tr\!\left(\widehat f_\lambda^\dagger
               \widehat q_\lambda\right).
  \label{eq:plancherel-appendix}
\end{equation}
Set $f=\cA$ and let $q=B$ be the convolution defined in
Eq.~\eqref{eq:convolution}.  The convolution theorem gives
$\widehat B_\lambda=\widehat A_\lambda K_\lambda$.  Substitution into
Eq.~\eqref{eq:plancherel-appendix} immediately yields
Eq.~\eqref{eq:block-quadratic-form}.

The placement of inverses and whether matrices are stored transposed varies
between Fourier conventions.  The source code fixes one consistent choice:
function values are in lexicographic image-permutation order, the forward
transform contains $\rho(g)$ rather than $\rho(g^{-1})$, and the kernel
acts by the right convolution in Eq.~\eqref{eq:convolution}.  With those
choices, the block product is $\widehat A_\lambda K_\lambda$.

\section{Glossary}

\begin{description}[leftmargin=3.2cm,style=nextline]
  \item[Colour order]
    A permutation specifying the order used in a trace or DDM partial
    amplitude.
  \item[Gram matrix]
    The matrix of pairwise inner products between colour tensors.
  \item[Kernel]
    The one-argument function $k(g^{-1}h)$ that determines every Gram
    entry.
  \item[Convolution]
    A sum in which a kernel depends only on the relative group element.
  \item[Representation]
    Matrices that reproduce the multiplication law of a group.
  \item[Irrep]
    An irreducible representation, meaning an elementary symmetry sector.
  \item[Young diagram]
    A partition-shaped label for an irrep of a symmetric group.
  \item[Fourier block]
    The matrix coefficient associated with one irrep.
  \item[Plancherel identity]
    The equality that converts an inner product of group functions into a
    weighted sum of inner products of their Fourier blocks.
  \item[DDM basis]
    A colour decomposition with two fixed endpoint gluons and a permutation
    of the gluons between them.
\end{description}

\end{document}
